Adaptation helps legged robots survive distribution shift, but adaptation itself consumes battery, latency, and onboard compute. This paper studies a deployment question: when battery is limited, should a quadruped adapt continuously, or should adaptation be scheduled as a resource-budgeted action? We answer that question with an executable stochastic simulator shipped in this repository. The simulator models latent shift, battery depletion, adaptation latency, and compute energy across terrain and payload disturbances at both $25\%$ and $80\%$ battery, with $1024$ episodes per regime-policy pair. Compared with always-on adaptation, the battery-gated policy improves mission return from $0.54$ to $0.56$ in low-battery terrain shift and from $0.60$ to $0.63$ in low-battery payload shift while reducing adaptation-compute energy from $0.41$ to $0.15$ Wh and from $0.42$ to $0.07$ Wh, respectively. In high-battery regimes, the learned threshold saturates to effectively always-on behavior and matches its $0.99$ mission return. A cost sweep shows that gating is strongest at moderate adaptation cost, while both policies degrade once adaptation becomes extremely expensive, exposing a real limitation of a single fixed threshold. The result is still modest: this is a lightweight simulator, not a real robot benchmark. But it is now a real computation-backed study with code, episode-level data, generated figures, and a fully reproducible paper package.
